\chapter{Health Probes and Resource Management}
\label{ch:probes}

This chapter is the operational heart of Part~V: how the kubelet decides
a pod is alive and ready, and how the scheduler and kernel share two
cores among a dozen JVMs. Nearly every line was learned by watching a
real pod misbehave.

\section{The three probes}

\begin{description}
  \item[startupProbe] ``still booting --- be patient.'' While it fails,
  the other probes are suspended. Exhaust its budget and the container
  is restarted.
  \item[readinessProbe] ``may traffic come \emph{right now}?'' Failing
  pulls the pod out of Service endpoints --- no restart, no drama, just
  no new requests until it passes again.
  \item[livenessProbe] ``is the process beyond saving?'' Failing
  \emph{restarts} the container. The dangerous one.
\end{description}

\begin{figure}[h]
\centering
\begin{tikzpicture}[node distance=4mm and 12mm]
  \node[box] (boot) {starting\\\scriptsize startup probe failing};
  \node[boxfill, right=of boot] (ready) {Ready\\\scriptsize in Service endpoints};
  \node[box, right=of ready] (notready) {not Ready\\\scriptsize removed from endpoints};
  \draw[flow] (boot) -- node[lbl,above]{startup passes} (ready);
  \draw[flow, bend left=18] (ready) to node[lbl,above]{readiness fails} (notready);
  \draw[flow, bend left=18] (notready) to node[lbl,below]{readiness passes} (ready);
  \draw[flow, bend right=25] (boot.south) to node[lbl,below]{budget exhausted: restart} ++(0,-0.75);
\end{tikzpicture}
\caption{The pod's life as the kubelet sees it.}
\end{figure}

From \code{course-service.yaml}:

\begin{lstlisting}[language=yaml]
startupProbe:
  httpGet: { path: /v3/api-docs, port: http }   (*@\co{1}@*)
  periodSeconds: 10
  failureThreshold: 12                          (*@\co{2}@*)
readinessProbe:
  httpGet: { path: /v3/api-docs, port: http }
  periodSeconds: 5
  failureThreshold: 3
# no livenessProbe                              (*@\co{3}@*)
\end{lstlisting}

\begin{description}
  \item[\co{1}] The probe-path decision. \code{/actuator/health} is the
  purpose-built endpoint --- but course-service and auth-service do not
  permit it through Spring Security, so it answers 401 and the kubelet
  counts that as failure. The one unauthenticated path that only answers
  200 \emph{after the full context is up} --- Flyway migrated, beans
  built --- is springdoc's \code{/v3/api-docs}. Dictionary-service and
  the gateway \emph{do} permit actuator and probe it properly; and
  notification-service has no Spring Security at all, so its health
  endpoint is open by default --- three services, three answers to the
  same question. Rule:
  a probe path must be unauthenticated \emph{and} a true readiness
  signal; a \code{tcpSocket} check passing the moment Tomcat binds is
  neither.
  \item[\co{2}] $12 \times 10$\,s = two minutes of boot grace --- sized
  to a JVM cold-starting on two contended cores.
  \item[\co{3}] Deliberate absence. Liveness kills; the only thing it
  would catch beyond readiness is a truly wedged JVM --- rare --- while a
  long GC pause on a loaded 2-core box could fail the probe and kill a
  \emph{healthy} pod. Add liveness when you have observed the hang it
  is for.
\end{description}

\begin{pitfall}[Startup probe failed: statuscode 404 --- or 403]
Symptom (milestone~6, builds \#6): notification-service logs show
Kafka partitions assigned and a fully started application, yet the
pod never goes Ready; the probe reports \code{404}. dictionary-service,
same story, but \code{403}. Cause: the config-server files
\emph{expose} \code{health,info,metrics,prometheus} for every service
--- but only config-server, discovery-server and api-gateway had
\code{spring-boot-starter-actuator} on the classpath. Exposure
configuration for an absent library is silently inert; the URL is
just a 404. The 403 variant is Spring Security's contribution: a 404
forwards to \code{/error}, \code{/error} is not in the permitAll list,
and a stateless chain with no entry point answers 403 --- so the
symptom points at authorization while the cause is a missing jar.
Fix: the starter in all four service POMs. Lesson: a probe path must
be verified with \code{curl} from inside a pod \emph{before} it is
trusted in a manifest --- the book's earlier claim that
dictionary-service ``probes health properly'' was written from its
SecurityConfig, not from a live request.
\end{pitfall}

\begin{pitfall}[the restart loop that creates itself]
Symptom: a service restarts every few minutes under load, each restart
making the load worse. Cause: an aggressive liveness probe timing out
during GC pauses or slow downstream calls --- the probe \emph{causes}
the outage it was meant to detect. This failure mode is why
``no liveness probe'' is a defensible default and why liveness, when
added, should check only ``is the event loop responding,'' never
downstream dependencies.
\end{pitfall}

\section{Requests, limits, and the JVM}

\begin{lstlisting}[language=yaml]
resources:
  requests: { memory: 512Mi, cpu: 200m }   (*@\co{1}@*)
  limits: { memory: 768Mi }                (*@\co{2}@*)
env:
  - { name: JAVA_TOOL_OPTIONS, value: "-XX:MaxRAMPercentage=75.0" } (*@\co{3}@*)
\end{lstlisting}

\begin{description}
  \item[\co{1}] \emph{Requests} are the scheduler's accounting: the pod
  is placed only where 512\,Mi and 0.2 CPU are unclaimed. Requests
  \emph{reserve}, they do not cap.
  \item[\co{2}] \emph{Limits} are kernel-enforced ceilings. Exceed the
  memory limit and the process is OOM-killed --- status
  \code{OOMKilled}, exit code 137, no Java exception, no log line. Note
  the asymmetry: a memory limit, but \emph{no CPU limit} --- next
  decision box.
  \item[\co{3}] The container-aware JVM sizes its heap as a fraction of
  the \emph{container's} memory limit: 75\% of 768\,Mi for heap, the
  rest for metaspace, threads, and off-heap. Without this, defaults
  leave either waste or an OOM-kill margin of zero.
\end{description}

\begin{decision}[why no CPU limit]
CPU limits throttle: a container over its quota is paused until the
next scheduling period --- and a starting JVM is exactly when a service
wants every cycle (class loading, JIT). Throttled startup turns 60
seconds into 200 and fails the startup probe: the limit manufactures
the failure. Unlike memory, CPU is compressible --- contention slows
things instead of killing them --- so the pattern used here is standard
advice: \textbf{always request CPU, rarely limit it}. Memory gets a
limit because the failure mode of unbounded memory is the kernel
killing \emph{someone}, and you want the boundary predictable.
\end{decision}

\begin{hardware}
The requests across all services sum to roughly what 19\,GB minus
Jenkins, k3s, and the datastores can hold --- deliberately. Requests are
promises to the scheduler; over-promising on a single node means pods
that never schedule (\code{Pending} forever), under-promising means
OOM roulette. On one node you are doing capacity planning by hand;
a cloud autoscaler turns the same arithmetic into node-pool sizing.
\end{hardware}

\begin{handson}[watch a startup through the probe's eyes]
\begin{lstlisting}[language=bash]
kubectl -n ts delete pod -l app=course-service
kubectl -n ts get pods -w
# Running 0/1  <- startup probe still failing; no traffic yet
# Running 1/1  <- /v3/api-docs answered; endpoints updated
kubectl -n ts describe pod -l app=course-service | tail -n 12
\end{lstlisting}
The gap between \code{0/1} and \code{1/1} is Flyway plus context
startup --- the window in which, probe-less, requests would have hit a
half-started service.
\end{handson}

\begin{quiz}
\begin{enumerate}
  \item For each probe: what does failing it do, and which is safe to
  be aggressive with?
  \item Why is \code{/v3/api-docs} an honest readiness signal for
  course-service while \code{tcpSocket: 8082} is a lie? What makes
  \code{/actuator/health} unusable there, and usable on
  dictionary-service?
  \item A pod shows \code{OOMKilled}, exit 137, clean logs. Explain the
  mechanism and why no Java exception exists. First two things to check?
  \item Requests vs.\ limits: which does the scheduler read, which does
  the kernel enforce, and why does this project limit memory but not
  CPU?
\end{enumerate}
\end{quiz}
